% GENERATED FILE. Edit canonical lessons, then run npm run generate:books.
% canonical-source-hash: fnv1a64:2949289880ce15c7
% canonical-lessons: GE-C20-entschuldigung

\chapter{Sorry}
\label{ch:ge-sorry}

This chapter is generated from the canonical micro-lessons used by Language Ladder.
Each section stays independently resumable and preserves the authored prerequisite order.

\section[Entschuldigung]{Entschuldigung — “sorry” / “excuse me”}

\label{lesson:GE-C20-entschuldigung}



\begin{quote}
\textbf{Warm-up.} {[}PAUSE 2s{]} You already know \textbf{bitte} does triple duty for \textquotedblleft{}please,\textquotedblright{} \textquotedblleft{}you're welcome,\textquotedblright{} and \textquotedblleft{}here you go.\textquotedblright{} German's word for \textbf{sorry} is just as hard-working: \textbf{Entschuldigung} covers both \textquotedblleft{}\textbf{excuse me}\textquotedblright{} and \textquotedblleft{}\textbf{I'm sorry}.\textquotedblright{}
\end{quote}

\begin{cousinweb}
\begin{itemize}
\raggedright
  \item \textbf{ent-} = a \textquotedblleft{}removal / undoing\textquotedblright{} prefix — the same shape as in \textbf{ent}decken (\textquotedblleft{}un-cover\textquotedblright{} $\to$ discover) or \textbf{ent}fernen (\textquotedblleft{}un-distance\textquotedblright{} $\to$ remove).
  \item \textbf{Schuld} = \textquotedblleft{}\textbf{guilt, fault, blame}\textquotedblright{} (also, separately, \textquotedblleft{}debt\textquotedblright{} — a \emph{Schuld} can be money owed or wrong done).
  \item \textbf{-igung} = a noun-forming suffix, roughly \textquotedblleft{}-ing/-ation.\textquotedblright{}
\end{itemize}

So \textbf{Entschuldigung} is literally \textquotedblleft{}\textbf{an un-guilting}\textquotedblright{} — the act of removing blame. Saying it hands the guilt back: \emph{I release you from any fault} (when getting someone's attention or brushing past) or \emph{I ask you to release me from mine} (when apologizing).

\textbf{Schuld} has a striking cousin: it traces to the same ancient Germanic root as English \textbf{shall} and \textbf{should} — a verb that originally meant \textbf{\textquotedblleft{}to owe.\textquotedblright{}} \emph{Schuld} (what is owed, guilt) and \emph{shall} (what one is obliged to do) are, at root, the same idea of an obligation outstanding.
\end{cousinweb}

\begin{culture}
For \textbf{deeper regret} — real sorrow, not a bump in the hallway — German often reaches for \textbf{es tut mir leid}, literally \textquotedblleft{}\textbf{it does to-me sorrow}\textquotedblright{}:

\begin{itemize}
\raggedright
  \item \textbf{es} = \textquotedblleft{}it,\textquotedblright{} \textbf{tut} = \textquotedblleft{}does\textquotedblright{} (from \textbf{tun}, \textquotedblleft{}to do\textquotedblright{}), \textbf{mir} = \textquotedblleft{}to me,\textquotedblright{} \textbf{Leid} = \textquotedblleft{}sorrow, suffering\textquotedblright{} (a cousin of English \textbf{loath}, both from a root meaning \textquotedblleft{}hateful, painful\textquotedblright{}).
\end{itemize}

Use \textbf{Entschuldigung} for the everyday sorry/excuse-me (stepping on a foot, interrupting), and \textbf{es tut mir leid} when you mean it more — condolences, a real regret. The two overlap; native speakers move between them by feel.
\end{culture}

\subsection*{Guided Practice}
{[}PAUSE 1s{]}

\begin{itemize}
\raggedright
  \item {[}YOU SAY: \textquotedblleft{}Schuld\textquotedblright{} — guilt/fault — then \textquotedblleft{}Entschuldigung,\textquotedblright{} sorry / excuse me{]}
  \item {[}YOU SAY: the deeper phrase — \textquotedblleft{}es tut mir leid,\textquotedblright{} it pains me / I'm so sorry{]}
\end{itemize}

\begin{tcolorbox}[breakable,colback=teal!4,colframe=teal!35!black,title={Wrap-up recall}]
{[}PAUSE 3s{]} What does \textbf{Entschuldigung} literally mean, piece by piece? (\textbf{\textquotedblleft{}Un-guilting\textquotedblright{}} — \emph{ent-} \textquotedblleft{}away\textquotedblright{} + \emph{Schuld} \textquotedblleft{}guilt/fault.\textquotedblright{}) What English word shares Schuld's ancient root? (\textbf{Shall/should} — both from a root meaning \textquotedblleft{}to owe.\textquotedblright{}) What's the deeper phrase for real regret? (\textbf{Es tut mir leid}, \textquotedblleft{}it does to me sorrow.\textquotedblright{})
\end{tcolorbox}
